% Flujo General del Sistema - Versión Técnica Mejorada
% Refleja la implementación real con procesamiento por fases

\section{Flujo General del Sistema}

El sistema SGM-Payroll implementa un flujo de procesamiento secuencial por fases, diseñado específicamente para optimizar el manejo de grandes volúmenes de datos de nómina mientras mantiene la integridad y trazabilidad del proceso.

\subsection{Arquitectura del Flujo}

El procesamiento se organiza en un pipeline de 5 fases secuenciales, donde cada fase tiene responsabilidades específicas y utiliza estrategias de almacenamiento optimizadas:

\begin{enumerate}
    \item \textbf{Fase 1}: Análisis de Headers y Mapeo de Conceptos
    \item \textbf{Fase 2}: Procesamiento de Movimientos Talana  
    \item \textbf{Fase 3}: Integración de Archivos del Analista
    \item \textbf{Fase 4}: Mapeo y Sincronización de Novedades
    \item \textbf{Fase 5}: Verificación de Discrepancias y Consolidación
\end{enumerate}

\subsection{Flujo Detallado por Fases}

\subsubsection{Acceso y Gestión Inicial}

\begin{enumerate}
    \item \textbf{Autenticación}: Login con roles diferenciados (Analista, Supervisor, Gerente)
    \item \textbf{Asignación Automática}: El sistema asigna analista\_responsable automáticamente durante creación de cierres
    \item \textbf{Filtrado por Cliente}: Cada usuario ve solo los clientes y cierres asignados
    \item \textbf{Gestión de Estados}: Los cierres transicionan por estados: pendiente → en\_proceso → completado
\end{enumerate}

\subsubsection{Fase 1: Análisis de Headers y Mapeo}

\textbf{Objetivo}: Extraer conceptos del Libro de Remuneraciones y mapearlos sin procesar empleados aún.

\begin{itemize}
    \item \textbf{Input}: Libro de Remuneraciones Talana (Excel)
    \item \textbf{Procesamiento}:
    \begin{enumerate}
        \item Celery Task \texttt{analizar\_headers\_libro} extrae 55 conceptos detectados
        \item Filtrado automático de columnas de empleado (RUT, nombres, etc.)
        \item Mapeo automático usando tabla \texttt{MapeoConcepto} del cliente
        \item Identificación de conceptos pendientes de mapeo manual
    \end{enumerate}
    \item \textbf{Almacenamiento Redis}:
    \begin{verbatim}
    sgm:nomina:{cliente}:{cierre}:fase1_headers:analisis
    sgm:nomina:{cliente}:{cierre}:fase1_headers:mapeo_final
    sgm:nomina:{cliente}:{cierre}:estado:procesamiento
    \end{verbatim}
    \item \textbf{Optimización}: Solo 5KB de metadatos vs 365KB de datos completos (85\% reducción)
    \item \textbf{Output}: Mapeo consolidado listo para siguientes fases
\end{itemize}

\subsubsection{Fase 2: Movimientos Talana}

\textbf{Objetivo}: Procesar ausentismos, ingresos y finiquitos de empleados.

\begin{itemize}
    \item \textbf{Filtrado Inteligente}: Solo empleados que existen en Libro de Remuneraciones
    \item \textbf{Validaciones}: Empleados fuera del scope se excluyen automáticamente
    \item \textbf{Almacenamiento}: PostgreSQL para persistencia, cache selectivo en Redis
    \item \textbf{Estructura Redis}:
    \begin{verbatim}
    sgm:nomina:{cliente}:{cierre}:fase2_movimientos:talana
    \end{verbatim}
\end{itemize}

\subsubsection{Fase 3: Archivos del Analista}

\textbf{Objetivo}: Integrar archivos del analista aplicando filtrado por empleados existentes.

\begin{itemize}
    \item \textbf{Archivos}: Finiquitos, Ingresos, Ausentismos del analista
    \item \textbf{Validación Cruzada}: Solo empleados presentes en lista del Libro
    \item \textbf{Scope Control}: Datos fuera del proyecto se excluyen automáticamente
\end{itemize}

\subsubsection{Fase 4: Novedades y Mapeo de Headers}

\textbf{Objetivo}: Sincronizar términos entre archivos del analista y Talana.

\begin{itemize}
    \item \textbf{Mapeo de Headers}: Conceptos del analista ↔ conceptos del Libro
    \item \textbf{Lista de Empleados}: Nueva lista de empleados + mapeo de headers
    \item \textbf{Sincronización}: Normalización de terminología entre fuentes
\end{itemize}

\subsubsection{Fase 5: Verificación de Discrepancias}

\textbf{Objetivo}: Comparar datos entre fuentes y detectar inconsistencias.

\begin{itemize}
    \item \textbf{Comparaciones Realizadas}:
    \begin{enumerate}
        \item Movimientos (Talana) vs Datos de Analista (ingresos, finiquitos, ausentismos)
        \item Concepto-Valor (Libro Talana) vs Concepto-Valor (Novedades Analista)
        \item Empleados que pueden no coincidir entre fuentes
    \end{enumerate}
    \item \textbf{Cache Temporal}: Solo durante comparación activa
    \item \textbf{Cleanup Automático}: Redis se limpia post-verificación
    \item \textbf{Persistencia}: Discrepancias finales en PostgreSQL para auditoría
\end{itemize}

\subsection{Gestión de Estados y Transiciones}

El sistema implementa una máquina de estados sofisticada para tracking del progreso:

\begin{verbatim}
Estados del Cierre:
pendiente → headers_mapeados → listo_procesar → fase1_completa 
→ movimientos_cargados → analista_cargado → verificacion_completada

Estados de Error:
mapeo_requerido → error → reintentar
\end{verbatim}

\subsection{Optimizaciones de Performance}

\subsubsection{Procesamiento Asíncrono}
\begin{itemize}
    \item \textbf{Celery Chains}: Tareas encadenadas con passing de resultados optimizado
    \item \textbf{Retry Logic}: Reintentos automáticos con backoff exponencial
    \item \textbf{Circuit Breaker}: Prevención de cascading failures
\end{itemize}

\subsubsection{Gestión de Memoria}
\begin{itemize}
    \item \textbf{Cache Selectivo}: Solo datos necesarios para fase activa
    \item \textbf{TTL Diferenciado}: 1h para verificación, 24h para mapeos críticos
    \item \textbf{Cleanup Automático}: Liberación de recursos post-proceso
\end{itemize}

\subsubsection{Escalabilidad por Fases}
\begin{itemize}
    \item \textbf{Filtrado Automático}: Reduce volumen de datos por fase
    \item \textbf{Procesamiento Incremental}: Sin duplicación de datos entre fases
    \item \textbf{Particionamiento}: Por cliente y período para distribución eficiente
\end{itemize}

\subsection{Validaciones y Control de Calidad}

\subsubsection{Validaciones por Fase}
\begin{itemize}
    \item \textbf{Fase 1}: Formato Excel, estructura de headers, RUTs válidos
    \item \textbf{Fase 2-4}: Empleados existentes, consistencia de fechas
    \item \textbf{Fase 5}: Integridad referencial, rangos de valores
\end{itemize}

\subsubsection{Manejo de Errores}
\begin{itemize}
    \item \textbf{Rollback por Fases}: Reversión granular sin afectar fases completadas
    \item \textbf{Logging Estructurado}: Trazabilidad completa de errores y decisiones
    \item \textbf{Recovery State}: Recuperación desde cache en caso de fallos
\end{itemize}

\subsection{Generación de KPIs y Reporting}

Una vez completadas todas las fases, el sistema calcula automáticamente:

\begin{itemize}
    \item \textbf{KPIs Operacionales}: Dotación, haberes, descuentos, aportes patronales
    \item \textbf{Métricas de Performance}: Tiempo por fase, empleados procesados, discrepancias detectadas
    \item \textbf{Análisis Comparativo}: Variaciones respecto al período anterior
    \item \textbf{Dashboards Ejecutivos}: Consolidación para gerentes y supervisores
\end{itemize}

Este enfoque por fases representa una innovación metodológica específica para el procesamiento de nóminas, donde la optimización de recursos y la organización secuencial son fundamentales para manejar la complejidad inherente del dominio.
